\begin{figure*}[t]
\centering
\resizebox{\textwidth}{!}{%
\begin{tikzpicture}[
  font=\footnotesize,
  stage/.style={draw=black!40, rounded corners=2pt, fill=#1, align=left, inner sep=4pt, minimum height=3.1cm, anchor=north west},
  exper/.style={draw=black!40, rounded corners=2pt, fill=#1, align=left, inner sep=4pt, text width=5.35cm, minimum height=2.55cm, anchor=north west},
  arr/.style={-{Stealth[length=2mm]}, thick, black!55},
]
\coordinate (bus) at (0,-3.5);
\coordinate (row2) at (0,-3.9);

\node[stage=blue!4, text width=3.6cm] (data) at (0,0)
  {\textbf{PaHaW recordings}\\[2pt]
   75 participants, 8 tasks,\\597 recordings\\[3pt]
   per sample: position, time, pen contact, pressure, altitude and azimuth};

\node[stage=orange!7, text width=4.0cm, right=5mm of data.north east, anchor=north west] (render)
  {\textbf{Pseudodynamic minimap}\\[2pt]
   \textbullet\ trace drawn only where the pen touches the tablet\\
   \textbullet\ each signal normalized within its own recording\\
   \textbullet\ azimuth centered on its circular mean\\
   \textbullet\ up to three signals mapped to RGB};

\node[stage=orange!7, text width=3.9cm, right=5mm of render.north east, anchor=north west] (enc)
  {\textbf{Nine encodings} (Table~\ref{tab:encoding})\\[3pt]
   \begin{tabular}{@{}c@{\hspace{8pt}}c@{}}
   \includegraphics[height=1.2cm]{method_static.png} & \includegraphics[height=1.2cm]{method_saz.png}\\
   \scriptsize 1: static trace & \scriptsize 8: $SAZ$
   \end{tabular}\\[2pt]
   2--5: one signal; 6--9: RGB triplets};

\node[stage=blue!4, text width=3.7cm, right=5mm of enc.north east, anchor=north west] (eval)
  {\textbf{Evaluation}\\[2pt]
   $5\times5$ participant-separated cross-validation with 3 inner folds\\[3pt]
   recording-level macro F1\\[1pt]
   participant-level macro F1 (majority vote over tasks)};

\node[exper=green!6] (e1) at (data.north west |- row2)
  {\textbf{Experiment 1} \hfill Sec.~\ref{sec:descriptor}--\ref{sec:classifiers}\\[2pt]
   all eight tasks\\[2pt]
   \textbullet\ each encoding described by LPQ texture and classified by SVM\\
   \textbullet\ leading encoding ($SAZ$): SVM compared with logistic regression, random forest and a CNN trained from scratch};

\node[exper=green!6, right=5mm of e1.north east, anchor=north west] (e2)
  {\textbf{Experiment 2} \hfill Sec.~\ref{sec:nested-followup}--\ref{sec:controls}\\[2pt]
   all eight tasks\\[2pt]
   \textbullet\ encoding, SVM and fusion rule chosen with training participants only\\
   \textbullet\ spatial controls: signals circularly shifted or permuted within each stroke};

\node[exper=green!6, right=5mm of e2.north east, anchor=north west] (e3)
  {\textbf{Experiment 3} \hfill Sec.~\ref{sec:competitive}\\[2pt]
   tasks T2--T4 (all eight as secondary)\\[2pt]
   \textbullet\ static trace + LPQ + SVM\\
   \textbullet\ minimaps + LPQ + SVM\\
   \textbullet\ minimaps + ResNet18 features\\
   \textbullet\ kinematic/pressure attributes + SVM};

\draw[arr] ([yshift=-1.55cm]data.north east) -- ([yshift=-1.55cm]render.north west);
\draw[arr] ([yshift=-1.55cm]render.north east) -- ([yshift=-1.55cm]enc.north west);
\draw[arr] ([yshift=-1.55cm]enc.north east) -- ([yshift=-1.55cm]eval.north west);

\draw[thick, black!55] (eval.south) -- (eval.south |- bus);
\draw[thick, black!55] (e1.north |- bus) -- (eval.south |- bus);
\foreach \e in {e1,e2,e3} \draw[arr] (\e.north |- bus) -- (\e.north);

\draw[arr, dashed, black!45] (data.north) -- ++(0,0.45)
  -| ([xshift=5mm,yshift=-2.05cm]e3.north east)
  node[pos=0.25, above, font=\scriptsize, text=black!70] {pen signals used directly, without an image}
  -- ([yshift=-2.05cm]e3.north east);
\end{tikzpicture}}
\caption{Overview of the method. Each PaHaW recording is rendered as a pseudodynamic minimap in one of nine encodings, and every experiment uses the same participant-separated cross-validation. Experiment~1 compares encodings and classifiers, Experiment~2 chooses the encoding with training participants only and tests whether the position of the signals along the stroke matters, and Experiment~3 compares minimaps with the static trace and with kinematic/pressure attributes computed directly from the pen signals (dashed arrow). Thumbnails show the spiral of one healthy participant as rendered.}
\label{fig:method}
\end{figure*}
